\documentclass[11pt]{article}
\usepackage[margin=1in]{geometry}
\usepackage{amsmath,amssymb,booktabs,hyperref}
\title{Heterogeneous Active Retrieval Across Structurally Different Evidence Substrates}
\author{Working Draft}\date{August 2026}
\begin{document}\maketitle
\begin{abstract}
FLARE and Self-RAG establish important forms of adaptive generation-time
retrieval \cite{flare,selfrag}. This paper does not claim that mechanism. It
tests whether structurally different information needs benefit from routing
among genuinely heterogeneous evidence substrates: relational execution,
lexical indexing, vector retrieval, web search, and optionally graph traversal.
The proposed router remains transparent and heuristic. The decisive controls
are matched-cost universal retrieval, broadcast, explicit model selection, and
oracle source routing. If oracle heterogeneous routing cannot beat a universal
retriever, later learned routing is not motivated.
\end{abstract}

\section{Architecture}
The runtime forms a typed information-need event
\[
h_t\rightarrow e_t=(\text{entity},\text{relation},\text{time},\text{source hints},\text{confidence}),
\]
then a transparent router chooses among
\[
\{\mathrm{DB},\mathrm{lexical},\mathrm{vector},\mathrm{KG},\mathrm{web}\}.
\]
Each event is logged through the factorization
\[
\textsc{trigger}\rightarrow\textsc{source}\rightarrow\textsc{query}
\rightarrow\textsc{retrieve/execute}\rightarrow\textsc{evidence\ use}.
\]
Component scores and oracle substitutions prevent routing failures from being
hidden inside final-answer accuracy.

\section{Minimum Retrieval Suite}
\begin{itemize}
\item relational/structured database;
\item lexical/document index;
\item vector database or semantic document retriever;
\item web/search source.
\end{itemize}
A knowledge graph is optional if it adds a genuinely distinct access pattern.
Multiple wrappers over one text retriever do not satisfy this requirement.

\section{Heuristic Routing}
Examples:
\begin{itemize}
\item structured attributes, identifiers, aggregates $\rightarrow$ DB;
\item explicit document names/exact terms $\rightarrow$ lexical;
\item semantic document questions $\rightarrow$ vector;
\item latest/current/recent external facts $\rightarrow$ web;
\item entity/relation neighborhoods $\rightarrow$ KG.
\end{itemize}
The benchmark makes these semantics necessary: an argmax over sales requires DB
aggregation, an exact policy clause favors lexical lookup, a semantic report
question favors vector retrieval, a current public fact favors web search, and a
neighborhood or path query favors a KG.
No learned router is allowed in the core paper.

\section{Composed Helpers}
Small deterministic preprocessing may resolve relative dates, normalize entity
IDs, or form database filters/aggregations. Bounded compositions such as date
resolver $\rightarrow$ web and entity resolver $\rightarrow$ DB are allowed,
but their dependency DAG and component costs must be logged. They may not become
a hidden general planner.

\section{Common Evidence IR}
All retrievers emit source type/ID, query, content/value, timestamp/freshness, relevance score where meaningful, support/contradiction/ambiguity status, and provenance.

\section{Tasks}
Use source-diverse workloads: structured business facts, exact document lookup, semantic document QA, recent/current facts, multi-source conflicts, no-retrieval cases, and cases where the best source becomes clear only mid-generation.

\section{Baselines}
\begin{enumerate}
\item LLM only.
\item Single universal retriever.
\item Upfront RAG across all sources.
\item FLARE-like confidence trigger plus universal retriever.
\item FLARE-like confidence trigger plus heuristic heterogeneous router.
\item Self-RAG-inspired adaptive retrieval, if feasible and clearly labeled as
an approximation unless canonical training is reproduced.
\item Explicit LLM tool/source selection.
\item Broadcast-to-all under matched cost and fanout.
\item Heuristic semantic reflex routing.
\item Oracle source router.
\item Oracle trigger plus source router.
\end{enumerate}

\section{Metrics}
Report final accuracy and unsupported claims, plus separate trigger P/R, source
top-$1$/top-$k$ and confusion matrix, query correctness, retrieval/execution P/R,
evidence-use rate, source fanout, conflicts, cost, generated/tool-selection
tokens, latency, dependency-DAG cost, and scaling with source count and factual
subgoals.

\section{Falsification}
The heterogeneous claim fails if a universal retriever matches under equal cost,
oracle routing adds little, matched-cost broadcasting dominates, source-specific
query semantics do not matter, explicit source selection is equally efficient,
or routing errors dominate specialization gains.

\section{Scope}
No learned routing, hidden-state trigger, backtracking, or PRA/native KV integration.

\begin{thebibliography}{9}
\bibitem{flare}
Z.~Jiang et al.
Active retrieval augmented generation.
\emph{EMNLP}, 2023.
\url{https://arxiv.org/abs/2305.06983}.

\bibitem{selfrag}
A.~Asai et al.
Self-RAG: Learning to retrieve, generate, and critique through self-reflection.
\emph{ICLR}, 2024.
\url{https://arxiv.org/abs/2310.11511}.
\end{thebibliography}

\end{document}
